\subsection{Operação do manipulador robótico}
\label{sec:estrategia_operacao_MR}

A operação do manipulador robótico é organizada em fases, de acordo com o tipo de aproximação realizado.
Na fase de aproximação longa, o manipulador permanece travado em uma configuração fixa, e o conjunto base--\gls{mr} é tratado como um corpo rígido.
O controle de posição relativa e de atitude atua apenas por meio dos atuadores da base física, sem movimentação das juntas do \gls{mr}.

Durante a aproximação curta, ou seja, momento em que o perseguidor entra na região de captura, o \gls{mr} é liberado para operação.
A estratégia consiste:
(i) no posicionamento inicial do manipulador, partindo da configuração travada para uma pose de trabalho;
(ii) alinhamento fino da ferramenta em relação ao eixo da  de acoplamento, ajustando a orientação e a posição do efetuador dentro da região de captura; e
(iii) aproximação final e inserção, com pequenos deslocamentos controlados até o contato e acoplamento mecânico.

Ao longo dessas fases, os movimentos do manipulador geram forças e torques internos que se manifestam como reações na base, podendo perturbar a atitude mesmo quando o controle de atitude está atuando.
A estratégia de operação do \gls{mr} leva isso em conta ao privilegiar trajetórias suaves, com variações moderadas de velocidade articular.